% =====================================================================
% APPENDIX: supplementary figures and declaration of AI/LLM use
% =====================================================================

\appendix

\section{Supplementary figures}\label{app:figures}

\begin{figure}[h]
    \centering
    \includegraphics[width=\linewidth]{a_fits.pdf}
    \caption{Training data for seed 42, Runge's function, and OLS fits of degree 3, 8 and 15. The last training point is at $x = 0.951$; beyond it the degree-8 fit rises to 0.655 at $x = 1$, where $f(1) = 0.038$.}
    \label{fig:a_fits}
\end{figure}

\begin{figure*}
    \centering
    \includegraphics[width=0.85\textwidth]{a_extrapolation.pdf}
    \caption{Test MSE of each of the 200 repetitions against how far the test points reach outside the interval covered by the training points, for degree 10 (left) and 20 (right). Grey: all test points inside the training interval. $n = 100$, $\sigma = 0.1$.}
    \label{fig:a_extrapolation}
\end{figure*}

\section{Use of AI/LLM tools}\label{app:llm}

\paragraph*{Tools used.} Claude (Opus 5.5), accessed through claude.ai including its code-execution workspace, September 2026. Claude was also used as a tutor for the theory; the understanding this gave is not declared separately.

\paragraph*{Text.} Contribution levels as defined in the course guidelines.

\begin{table}[h]
    \centering
    \begin{tabular}{lc}
        \toprule
        Section & Level \\
        \midrule
        Abstract & 3 \\
        Introduction & 3 \\
        Methods & 3 \\
        Results and discussion & 3 \\
        Conclusion & 3 \\
        Appendix A (figure captions) & 3 \\
        Appendix B (this declaration) & 3 \\
        \bottomrule
    \end{tabular}
\end{table}

\emph{Level 3 note.} The text was drafted by Claude from our dialogue about the methods and the results. For the results, we interpreted the figures and the logged output of the code ourselves first; Claude corrected our interpretation and drafted the text from it. We read, corrected and approved all sections, and statements we could not explain ourselves were simplified or removed. Every number in the text was checked against the logged output of the code in a separate pass by Claude.

\paragraph*{Code.} Contribution levels as defined in the course guidelines. Every function and class carries an \texttt{LLM-assisted} tag in its docstring that refers to the note at the top of its file.

\begin{table}[h]
    \centering
    \begin{tabular}{lcp{0.55\linewidth}}
        \toprule
        File & Level & Description \\
        \midrule
        \texttt{common.py} & 4 & Data, design matrix, scaling, OLS and Ridge solvers, scikit-learn Lasso, error measures, bootstrap, scikit-learn wrapper. Written by Claude from choices agreed with the authors. \\
        \texttt{optim.py} & 4 & Cost, gradients (analytic and JAX), update rules, gradient descent and SGD, following the course's reference implementation. Written by Claude. \\
        \texttt{main.py} & 4 & One function per project part, figures and logged output. Written by Claude from the experiments agreed with the authors. \\
        \texttt{tests.py} & 4 & Unit tests. Written by Claude; the checks were chosen together with the authors. \\
        \bottomrule
    \end{tabular}
\end{table}

The results were first produced by running the code in Claude's code-execution workspace and then reproduced by running it locally; the numbers in this report are from the local run. The code was tested against scikit-learn, JAX and the closed-form solutions (Sec.~\ref{sec:validation}), and the choices behind the experiments (seeds, repetitions, grids, what to compare) were discussed and agreed on before the code was written.
